\documentclass[12pt,a4paper]{report}

\usepackage[utf8]{inputenc}
\usepackage[T1]{fontenc}
\usepackage{geometry}
\usepackage{booktabs}
\usepackage{array}
\usepackage{hyperref}
\usepackage{enumitem}
\usepackage{titlesec}
\usepackage{pgfgantt}
\usepackage{xcolor}
\geometry{margin=2.5cm}
\usepackage{unicode-math}
\setmathfont{Stix Two Math}
\titleformat{\chapter}{\normalfont\huge\bfseries}{\chaptername\ \thechapter}{20pt}{\huge}

\begin{document}

%------------------------------------------------
% Title Page
%------------------------------------------------

\begin{titlepage}
\centering

{\Large Master of Science in Engineering \par}
\vspace{0.5cm}
{\Large Project Definition Report \par}

\vspace{2cm}

{\Huge\bfseries Design and Integration of a Planar Transformer and Charge-Pump--Based PFC for a High-Power-Density USB-PD AC--DC Power Supply \par}

\vspace{2cm}

Start date: 1 February 2026 \\
End date: 1 July 2026 \\
ECTS points: 30 \\

\vfill

Technical University of Denmark \\
DTU Electro \\
Department of Electrical and Photonics Engineering \\

\vspace{1cm}

Kongens Lyngby, 2026

\end{titlepage}

%------------------------------------------------
\chapter*{Summary}
\addcontentsline{toc}{chapter}{Summary}

This thesis investigates the design and integration of a planar transformer and a charge-pump--based power factor correction (PFC) stage for a compact, high-power-density USB Power Delivery (USB-PD) AC--DC power supply. The objective is to develop a converter capable of delivering \textbf{100 W continuous output power}, with a \textbf{minimum acceptable output of 65 W}, while achieving high efficiency, stable thermal performance, and high power factor.

The project combines high-frequency planar magnetic design, wide-bandgap semiconductor devices, and alternative PFC architectures to evaluate whether charge-pump--based PFC can provide competitive power density compared to conventional boost PFC approaches.

The work includes analytical design, component selection, PCB layout for planar magnetics and high-frequency switching, prototype fabrication, and comprehensive experimental validation including efficiency, power factor, transient response, and thermal performance characterization.

\tableofcontents
\clearpage

%------------------------------------------------
\chapter{Project Description}

\section{Background}

USB Power Delivery has become a universal high-power interface supporting programmable voltages up to 20 V and currents up to 5 A. The demand for compact adapters capable of powering laptops and embedded systems imposes strict constraints on:

\begin{itemize}
    \item Power density
    \item Efficiency across multiple PDOs
    \item Thermal performance
    \item Power factor and THD
\end{itemize}

Conventional AC--DC architectures typically employ a boost PFC stage followed by an isolated DC--DC converter using wound magnetic components. These solutions are often limited in volumetric efficiency and thermal integration.

Planar transformers enable high-frequency operation, reduced height, improved thermal spreading, and manufacturability via PCB integration. Charge-pump--based PFC topologies may offer an alternative to classical boost PFC by reducing magnetic component size and potentially increasing integration density.

This thesis targets a USB-PD supply capable of \textbf{100 W continuous output power}, with a \textbf{minimum guaranteed output of 65 W}, under defined thermal and operational constraints.

\section{Prior Work}

Planar magnetics have been extensively investigated for high-frequency power conversion in telecom and server applications, where low profile and high efficiency are critical.

Boost PFC remains the dominant architecture for single-phase AC--DC conversion. Alternative approaches, including charge-pump and switched-capacitor techniques, have been proposed to reduce magnetic size and improve integration at moderate power levels.

Wide-bandgap devices such as GaN and SiC enable higher switching frequencies and improved efficiency, facilitating further size reduction of passive components. However, integrating planar magnetics with a charge-pump--based PFC in a USB-PD compliant system presents a multi-domain design challenge requiring careful analytical design and experimental validation.

\section{Problem Statement}

How can a planar transformer and a charge-pump--based PFC stage be analytically designed and experimentally validated to deliver \textbf{100 W (≥65 W minimum)} in a compact USB-PD AC--DC power supply while maximizing efficiency, power density, thermal stability, and power quality?

\subsection*{Sub-goals}

\begin{itemize}
    \item Define electrical and mechanical requirements including:
    \begin{itemize}
        \item USB-PD PDO profiles
        \item Efficiency targets
        \item Thermal constraints
        \item Power-factor requirements
    \end{itemize}

    \item Analyze and compare charge-pump PFC topologies against conventional boost PFC with respect to:
    \begin{itemize}
        \item Efficiency
        \item Component stress
        \item Power density
        \item Control complexity
    \end{itemize}

    \item Integrate and optimize a planar transformer considering:
    \begin{itemize}
        \item High-frequency operation
        \item Controlled leakage inductance
        \item Winding AC resistance minimization
        \item Core loss estimation
        \item Manufacturability and low-profile constraints
    \end{itemize}

    \item Develop the complete power stage architecture including:
    \begin{itemize}
        \item PFC stage
        \item LLC stage
        \item Rectification
        \item Control strategy
        \item USB-PD interface behavior
    \end{itemize}

    \item Fabricate and assemble a prototype including:
    \begin{itemize}
        \item PCB-integrated planar magnetics
        \item High-frequency layout practices
        \item Safety clearances and isolation
        \item Integrated thermal management
    \end{itemize}

    \item Experimentally characterize:
    \begin{itemize}
        \item Power factor and THD
        \item Efficiency across PDOs
        \item Transient response
        \item Thermal distribution and stability
        \item Reliability under extended load testing
    \end{itemize}

    \item Critically evaluate trade-offs and propose improvements in:
    \begin{itemize}
        \item Planar magnetic integration
        \item Wide-bandgap stage optimization
        \item Manufacturability at scale
        \item Future USB-PD evolution
    \end{itemize}
\end{itemize}

\subsection*{Success Criteria}

\begin{itemize}
    \item ≥ 65 W continuous output demonstrated in hardware.
    \item 100 W continuous operation achieved under nominal conditions.
    \item Stable operation across selected USB-PD PDOs verified.
    \item Power factor, efficiency, transient response, and thermal performance experimentally validated and documented.
\end{itemize}

%------------------------------------------------
\chapter{Intended Learning Objectives}

\begin{itemize}
    \item Demonstrate the ability to define and structure a high-frequency AC--DC power conversion system integrating magnetics, switching devices, control, and thermal management.
    \item Apply analytical design methods for planar magnetics and high-frequency power stages.
    \item Design planar magnetics optimized for electrical performance and manufacturability.
    \item Integrate wide-bandgap devices into a compact high-power converter.
    \item Perform rigorous experimental validation of converter performance.
    \item Critically assess efficiency, density, and thermal trade-offs.
    \item Communicate technical findings clearly in written and oral form.
    \item Prove the charge pump based PFC circuit can be used for usb pd applications. 
    \item Test the power delivery over all the usb pd standard and pps functions. 
\end{itemize}

%------------------------------------------------
\chapter{Plan}

\section{Activities}

\begin{table}[h]
\centering
\begin{tabular}{lll}
\toprule
Phase & Description & Period \\
\midrule
1 & Project definition and requirement analysis & W5--W7 \\
2 & Literature review and topology selection & W6--W9 \\
3 & Planar transformer design and optimization & W8--W13 \\
4 & Power stage architecture development & W10--W15 \\
5 & PCB design and layout & W14--W18 \\
6 & Prototype fabrication and assembly & W18--W20 \\
7 & Experimental characterization and refinement & W20--W24 \\
8 & Thesis writing and documentation & W16--W26 \\
\bottomrule
\end{tabular}
\end{table}

\section{Milestones}

\begin{table}[h]
\centering
\begin{tabular}{llll}
\toprule
Title & Description & Type & Week \\
\midrule
Project definition submitted & Approved project definition & TH & 7 \\
Topology finalized & Architecture frozen & D & 10 \\
Transformer design freeze & Magnetics finalized & D & 13 \\
PCB fabrication & Design sent to manufacturer & D & 18 \\
Prototype assembled & Hardware ready for test & AC & 20 \\
65 W demo achieved & Continuous output validated & D & 22 \\
100 W demo achieved & Nominal target validated & D & 24 \\
Thesis submission & Final report & TH & 26 \\
\bottomrule
\end{tabular}
\end{table}

\section{Risk Analysis}

\begin{table}[h]
\centering
\begin{tabular}{p{3cm} p{7cm} c}
\toprule
Risk & Description & Risk (1--5) \\
\midrule
Transformer loss underestimation & Excessive copper or core loss & 4 \\
Charge-pump PFC instability & Resonance or control issues & 3 \\
Thermal hotspot formation & Local overheating in planar windings & 3 \\
Component availability & WBG devices or cores delayed & 2 \\
\bottomrule
\end{tabular}
\end{table}


\begin{figure}[h]
\centering
\begin{ganttchart}[
    hgrid,
    vgrid,
    time slot format=isodate,
    x unit=0.45cm,
    y unit chart=0.55cm,
    y unit title=0.75cm,
    title height=1,
    bar height=0.6,
    bar/.append style={fill=gray!30},
    bar progress label font=\scriptsize,
    bar label font=\scriptsize,
    group label font=\bfseries\scriptsize,
    milestone label font=\scriptsize,
    milestone/.append style={fill=black},
    progress label text={\scriptsize\pgfmathprintnumber[fixed,precision=0]{#1}\%},
    link/.style={-latex, thick}
]{2026-01-26}{2026-06-29}

% Weeks 5..26 correspond roughly to dates starting late Jan; labels below are what matters
\gantttitle{Masters project (Weeks 5--26)}{155} \\
\gantttitlelist{5,...,26}{7} \\

% -----------------------------
% Project definition / planning
% -----------------------------
\ganttgroup{Project definition and planning}{2026-01-26}{2026-02-15} \\
\ganttbar[progress=0]{Project description / scope definition}{2026-01-26}{2026-02-15} \\
\ganttbar[progress=0]{Requirements (USB-PD/PPS, power targets, constraints)}{2026-01-26}{2026-02-08} \\
\ganttbar[progress=0]{Literature review: charge-pump PFC vs boost PFC}{2026-02-02}{2026-02-22} \\

% -----------------------------
% Architecture + parts selection
% -----------------------------
\ganttgroup{Architecture and key component selection}{2026-02-09}{2026-03-08} \\
\ganttbar[progress=0]{Architecture freeze (PFC + isolation + control approach)}{2026-02-09}{2026-02-22} \\
\ganttbar[progress=0]{Planar transformer selection + electrical characterization}{2026-02-16}{2026-03-08} \\
\ganttbar[progress=0]{WBG device + gate driver selection}{2026-02-16}{2026-03-01} \\
\ganttbar[progress=0]{USB-PD controller selection + PDO/PPS plan}{2026-02-23}{2026-03-08} \\

% -----------------------------
% Hardware implementation
% -----------------------------
\ganttgroup{Hardware implementation}{2026-03-02}{2026-04-19} \\
\ganttbar[progress=0]{Charge-pump PFC implementation (schematic + BOM)}{2026-03-02}{2026-03-22} \\
\ganttbar[progress=0]{Isolation / DC-DC stage implementation (schematic + BOM)}{2026-03-09}{2026-03-29} \\
\ganttbar[progress=0]{Create BOM + schematics review (design checks)}{2026-03-16}{2026-03-29} \\
\ganttbar[progress=0]{PCB layout: power stage + planar magnetics integration}{2026-03-23}{2026-04-12} \\
\ganttbar[progress=0]{DFM + ordering (PCB + parts)}{2026-04-06}{2026-04-19} \\

% -----------------------------
% Bring-up + integration
% -----------------------------
\ganttgroup{Bring-up and integration}{2026-04-20}{2026-05-24} \\
\ganttbar[progress=0]{Assembly + initial power-up}{2026-04-20}{2026-05-03} \\
\ganttbar[progress=0]{Closed-loop bring-up + stability checks}{2026-04-27}{2026-05-10} \\
\ganttbar[progress=0]{USB-PD negotiation bring-up (PDOs)}{2026-05-04}{2026-05-17} \\
\ganttbar[progress=0]{PPS mode validation (programmable V/I)}{2026-05-11}{2026-05-24} \\

% -----------------------------
% Measurement campaign
% -----------------------------
\ganttgroup{Experimental validation}{2026-05-18}{2026-06-22} \\
\ganttbar[progress=0]{75 W continuous validation}{2026-05-18}{2026-05-31} \\
\ganttbar[progress=0]{100 W continuous validation}{2026-05-25}{2026-06-07} \\
\ganttbar[progress=0]{Efficiency vs PDO/PPS sweep}{2026-05-25}{2026-06-14} \\
\ganttbar[progress=0]{PF + THD characterization vs load}{2026-06-01}{2026-06-14} \\
\ganttbar[progress=0]{Transient response testing}{2026-06-01}{2026-06-14} \\
\ganttbar[progress=0]{Thermal mapping + stability testing}{2026-06-01}{2026-06-22} \\
\ganttbar[progress=0]{Extended stress / durability testing}{2026-06-08}{2026-06-22} \\

% -----------------------------
% Thesis writing
% -----------------------------
\ganttgroup{Thesis writing}{2026-03-30}{2026-06-29} \\
\ganttbar[progress=0]{Writing: Introduction}{2026-03-30}{2026-04-12} \\
\ganttbar[progress=0]{Writing: Related work}{2026-04-06}{2026-04-19} \\
\ganttbar[progress=0]{Writing: Method / implementation}{2026-04-20}{2026-05-17} \\
\ganttbar[progress=0]{Writing: Results}{2026-05-18}{2026-06-14} \\
\ganttbar[progress=0]{Writing: Discussion}{2026-06-01}{2026-06-22} \\
\ganttbar[progress=0]{Writing: Conclusion + abstract}{2026-06-15}{2026-06-29} \\

% -----------------------------
% Buffer + hand-in
% -----------------------------
\ganttgroup{Buffer and hand-in}{2026-06-15}{2026-06-29} \\
\ganttbar[progress=0]{Buffer time}{2026-06-15}{2026-06-22} \\
\ganttbar[progress=0]{Corrections + final hand-in}{2026-06-22}{2026-06-29} \\

% Milestones
\ganttnewline
\ganttmilestone{Architecture frozen}{2026-02-22}
\ganttmilestone{PCB ordered}{2026-04-19}
\ganttmilestone{Prototype operational}{2026-05-03}
\ganttmilestone{75 W demo}{2026-05-31}
\ganttmilestone{100 W demo}{2026-06-07}
\ganttmilestone{Thesis hand-in}{2026-06-29}

\end{ganttchart}
\caption{Gantt chart for the MSc. thesis project (Weeks 5--26).}
\end{figure}
\end{document}
